\chapter{Inguinal Hernia Repair}
\section*{What Is This Test or Procedure?}
Inguinal hernia repair returns protruding tissue to the abdomen and reinforces the weakened groin wall, often with mesh.
\section*{Why This Test or Procedure Is Ordered}
It is performed for a painful, enlarging, irreducible, obstructed, or strangulated groin hernia and is one of the most common general-surgery operations.
\section*{How This Test or Procedure Is Performed}
Open repair uses a groin incision; laparoscopic or robotic repair uses small abdominal incisions. The hernia sac is reduced and the defect is closed or reinforced.
\section*{Risks and Recovery}
Recurrence, chronic groin pain, infection, bleeding, urinary retention, bowel or bladder injury, and testicular complications are possible. Many uncomplicated patients return home the same day.
\section*{References}
\begin{enumerate}\item MedlinePlus. Hernia Repair. U.S. National Library of Medicine; 2025.\item Current Medical Diagnosis and Treatment. McGraw-Hill; 2025.\end{enumerate}
\clearpage
